\documentclass[11pt]{article}
\textwidth 15.9cm
\textheight 23. cm
\addtolength{\oddsidemargin}{-11.5mm}
\addtolength{\evensidemargin}{-11.5mm}
\addtolength{\topmargin}{-19.0mm}

%\usepackage{amsmath}
\usepackage[ansinew]{inputenc}
\usepackage[bbgreekl]{mathbbol}
\usepackage{amsmath,amssymb}
\usepackage{graphicx}
\usepackage{hyperref}
\usepackage{url}
\usepackage[numbers,square]{natbib}
\usepackage{multirow}
\usepackage{bbm}
\usepackage{authblk}
\usepackage{subcaption}
\renewcommand{\Affilfont}{\footnotesize}
%\renewcommand{\Authfont}{\normalsize}

\usepackage{listings}
%\usepackage{stmaryrd}
\usepackage{comment}
%\usepackage{showlabels}
\usepackage{soul}
%\usepackage[dvipsnames]{xcolor}

%%%%%%%% MACROS
\include{GempicLatexMacros}

\begin{document}

\section{Examples}

This page contains a list of models implemented in GEMPICX. The source code together with sample input and diagnostic files can be found in the \verb|Examples| folder.
The provided Python files to analyze the simulation results can either be used directly or translated into Jupyter notebooks using jupytext via
\begin{lstlisting}
    jupytext --to ipynb Weibel.py 
\end{lstlisting}

\subsection{Vlasov-Maxwell}

The Vlasov-Maxwell system describes plasma particles in their self-consistent electromagnetic fields. The non-relativistic Vlasov equation for a particle species $s$ of charge $q_s$ and mass $m_s$ governs the evolution of a distribution function $f_s$ on phase space $(\bx, \bv) \in \Omega \times \mathbb{R}^3$ and is given by
\begin{equation}
    \label{eq:vlasov}
    \dt f_s(t,\bx,\bv) + \bv \cdot \nabla_\bx f_s(t,\bx,\bv) + \frac{q_s}{m_s} \left(\bs E(t,\bx) + \bv \times \bs B(t,\bx)\right) \cdot \nabla_\bv f_s(t,\bx,\bv) = 0.
\end{equation}

It is coupled to Maxwell's equations
\begin{equation}
    \label{eq:maxwell_example}
    \begin{aligned}
        \dt \bs D  & = \curl \bs H - \bs J, \\
        \dt \bs B  & = - \curl \bs E,       \\
        \Div \bs D & = \rho,                \\
        \Div \bs B & = 0,
    \end{aligned}
\end{equation}
written in four-field formulation. The two descriptions of each field are related through the constitutive relations, i.e., $\bs D = \varepsilon_0 \bs E$ and $\bs H = \frac{1}{\mu_0} \bs B$ with vacuum permittivity and permeability $\varepsilon_0$ and $\mu_0$, respectively.
The charge and current sources for Maxwell's equations are given by moments of the distribution function
\begin{equation}
    \label{eq:vlasov_moments}
    \rho(t,\bx) = \sum_s q_s \int f_s(t,\bx,\bv)\dd \bv, \qquad \bs J(t,\bx) = \sum_s q_s \int \bv f_s(t,\bx,\bv) \dd \bv.
\end{equation}
Under suitable boundary conditions, the Vlasov-Maxwell system conserves the energy
\begin{equation}
    \cH = \frac{1}{2} \int \bs E \cdot \bs D \dd\bx + \frac{1}{2} \int \bs B \cdot \bs H \dd\bx + \sum_s \frac{1}{2} m_s \int |\bv|^2 f_s \dd\bx \dd\bv.
\end{equation}

We apply a structure-preserving mimetic finite difference PIC discretization and obtain a discrete noncanonical Hamiltonian system. The resulting numerical scheme and details on the implementation are given in \ref{sec:splitting_vlasovmaxwell}.

As a simulation example, we demonstrate the Weibel instability, which includes a genuine magnetic effect driven by a temperature anisotropy. \ref{fig:vmaxweibelfourier} and \ref{fig:vmaxweibelenergy} are created by running the 3D simulation of the Vlasov-Maxwell example using the input file \verb|Examples/VlasovMaxwell/weibel.input|.

\begin{figure}
    \centering
    \includegraphics[width=0.7\linewidth]{../_static/graphics/Examples/VlasovMaxwell_WeibelFourierComp.png}
    \caption{Real and imaginary parts for the fundamental mode of the Weibel instability.}
    \label{fig:vmaxweibelfourier}
\end{figure}

\begin{figure}
    \centering
    \includegraphics[width=0.7\linewidth]{../_static/graphics/Examples/VlasovMaxwell_WeibelEnergy.png}
    \caption{Energy composition and conserved quantities for the Weibel instability simulation.}
    \label{fig:vmaxweibelenergy}
\end{figure}


\subsection{Electrostatic}

\subsubsection{The Vlasov-Poisson model}
The Vlasov-Poisson system describes plasma particles in their self-consistent electrostatic fields. The non-relativistic Vlasov equation for a particle species $s$ of charge $q_s$ and mass $m_s$ governs the evolution of a distribution function $f_s$ on phase space $(\bx, \bv) \in \Omega \times \mathbb{R}^3$ and is given by
\begin{equation}
    \label{eq:vlasov_VP}
    \dt f_s(t,\bx,\bv) + \bv \cdot \nabla_\bx f_s(t,\bx,\bv) + \frac{q_s}{m_s} \left(\bs E(t,\bx) + \bv \times \bs B(t,\bx)\right) \cdot \nabla_\bv f_s(t,\bx,\bv) = 0.
\end{equation}
It is coupled to the Poisson equation
\begin{equation}
    \label{eq:Poisson_example}
    - \Delta \phi =\rho(t,\bx)
\end{equation}
with $\bs E = -\nabla\phi$.
The charge density is obtained from the distribution function by
\begin{equation}
    \rho(t,\bx) = \sum_s q_s \int f_s(t,\bx,\bv)\dd \bv, \qquad \bs J(t,\bx) = \sum_s q_s \int \bv f_s(t,\bx,\bv) \dd \bv.
\end{equation}
Under suitable boundary conditions, the Vlasov-Maxwell system conserves the energy
\begin{equation}
    \cH = \frac{1}{2} \int |\bs E|^2 \dd\bx + \sum_s \frac{1}{2} m_s \int |\bv|^2 f_s \dd\bx \dd\bv.
\end{equation}

As simulation examples, we have three classical electrostatic problems: Landau damping for Vlasov-Poisson (VP), the two-stream instability and the bump-on-tail instability. The corresponding input files are
\verb|Examples/Electrostatic/landauVP.input|, \verb|Examples/Electrostatic/twoStream.input| and \verb|Examples/Electrostatic/bumpOnTail.input|.

\begin{figure}
    \centering
    \includegraphics[width=\linewidth]{../_static/graphics/Examples/ElectrostaticVP_Landau.png}
    \caption{Landau damping: Evolution of Fourier modes of the electric field (left) and the energy (right).}
    \label{fig:landauVP}
\end{figure}

\begin{figure}
    \centering
    \includegraphics[width=\linewidth]{../_static/graphics/Examples/ElectrostaticVP_TSI.png}
    \caption{Two stream instability $v_0=2.4$: Evolution of Fourier modes of electric field and time history.}
    \label{fig:TwoStreamVP}
\end{figure}

\begin{figure}
    \centering
    \includegraphics[width=\linewidth]{../_static/graphics/Examples/ElectrostaticVP_TSIstable.png}
    \caption{Two stream instability (stable case) $v_0=1.3$: Evolution of Fourier modes of the electric field (left) and energy and momentum over time (right).}
    \label{fig:TwoStreamStableVP}
\end{figure}

\begin{figure}
    \centering
    \includegraphics[width=\linewidth]{../_static/graphics/Examples/ElectrostaticVP_BOT.png}
    \caption{Bump on Tail instability: Evolution of Fourier modes of the electric field (left) and energy and momentum over time (right).}
    \label{fig:BOTVP}
\end{figure}

\subsubsection{The quasineutral electrostatic Vlasov model}

The quasineutral electrostatic Vlasov model describes the evolution of plasma particles in a quasineutral plasma, \textit{i.e.} such that the total charge density is 0.  The non-relativistic Vlasov equation for an ion species of charge $q$ and mass $m$ governs the evolution of the ion distribution function $f$ on phase space $(\bx, \bv) \in \Omega \times \mathbb{R}^3$ and is given by
\begin{equation}
    \label{eq:vlasov_QNE}
    \dt f(t,\bx,\bv) + \bv \cdot \nabla_\bx f(t,\bx,\bv) + \frac{q}{m} \left(\bs E(t,\bx) + \bv \times \bs B_0(\bx)\right) \cdot \nabla_\bv f(t,\bx,\bv) = 0,
\end{equation}
where $ \mathbf{B}_0$ is a given external magnetic field.
On the other hand, the electrons are considered adiabatic and follow the Maxwell-Boltzmann distribution, which yields
$n_e = n_0 \exp(q_e \phi/(k_B T_e))$, where $k_B$ is the Boltzmann constant. Assuming $q_e=-q_i$, the quasineutrality condition becomes
$$n_e=n =\int f(t,\bx,\bv)\dd\bv = n_0 e^{\frac{q_e \phi}{k_B T_e}}.$$
Linearizing this expression around $\phi=0$, we obtain
$$ n = n_0 \left(1 + \frac{q_e \phi}{k_B T_e}\right).$$
This allows us to express the electric field with respect to the ion distribution function as
\begin{equation}
    \label{eq:electric_field_QN}
    \mathbf{E} = -\nabla\phi = -\frac{k_B T_e}{n_0 q_e} \nabla n = -\frac{k_B T_e}{n_0 q_e} \nabla \int f \dd\bv.
\end{equation}

This model is tested by checking the Laudau damping for ion acoustic waves. The input file for this example is \texttt{landauIAW.input}. The analytical dispersion relation is solved in the notebook \texttt{landauIAWdisp.py} in the dispersion relations directory and the diagnostics for this example are computed in the notebook \texttt{landauIAW.py} and shown in figure \ref{fig:landauIAW}.

\begin{figure}
    \centering
    \includegraphics[width=\linewidth]{../_static/graphics/Examples/ElectrostaticQN_LandauIAW.png}
    \caption{Landau damping of ion acoustic waves: Evolution of Fourier modes of the electric field (left) and energy and momentum over time (right).}
    \label{fig:landauIAW}
\end{figure}


\subsection{Cold Plasma}

The cold plasma model consists of Maxwell's equations coupled to a linearized fluid description for the electrons. The continuous system in four-field formulation is given by
\begin{align}
    \dt \bs D   & = \curl \bs H - \bs J_e,                               \\
    \dt \bs B   & = - \curl \bs E,                                       \\
    \dt \bs J_e & = \omega_{pe}^2 \bs E + \bs J_e \times \bs \omega_{ce}
\end{align}
where $\omega_{pe} = \sqrt{\frac{e^2 n}{\varepsilon_0 m_e}}$, $\bs \omega_{ce} = \frac{e \bs B_0}{m_e}$ denote the electron plasma and cyclotron frequency, respectively. $\bs B_0$ is a given stationary background magnetic field. The constitutive relations for the electric and magnetic field are given by
\begin{align}
    \bs D = \varepsilon_0 \bs E, \quad \bs H = \frac{1}{\mu_0} \bs B.
\end{align}
With appropriate boundary conditions, the system conserves the energy
\begin{align}
    \cH = \frac{1}{2} \int \bs E \cdot \bs D \dd \bx + \frac{1}{2} \int \bs B \cdot \bs H \dd \bx + \frac{1}{2} \int \frac{\bs J_e}{\omega_{pe}^2} \cdot \bs J_e \dd \bx
\end{align}

We apply a structure-preserving mimetic finite difference discretization and obtain a discrete Hamiltonian system. The system and resulting splitting time integration scheme is given in \ref{sec:splitting_coldplasma}.

\ref{fig:coldplasma_dispersion} shows the 1D dispersion relation of our implementation. The simulation uses the input file \verb|Examples/ColdPlasma/coldPlasmaDispersion.input|.

\begin{figure}
    \centering
    \includegraphics[width=0.7\linewidth]{../_static/graphics/Examples/ColdPlasma_Dispersion.png}
    \caption{1D cold plasma dispersion relation for othogonal magnetic field orientations. Analytical solutions are shown as dashed lines.}
    \label{fig:coldplasma_dispersion}
\end{figure}


\subsection{Linearized Vlasov-Maxwell}
The Vlasov-Maxwell equations \ref*{eq:vlasov}, \ref*{eq:maxwell_example} can be linearized around a stationary solution $f_0$, $\bs E_0$, $\bs B_0$. Assuming that $\bs E_0 \equiv 0$ and $f_0 = \frac{n_0(\bx)}{\left(\sqrt{2\pi} v_\text{th}\right)^3} e^{-\frac{\abs{\bv}^2}{2 v_\text{th}^2}}$ is a Maxwellian, plugging the expansions $f = f_0 + \epsilon f_1$, $\bs E = \bs E_1$, $\bs B = \bs B_0 + \epsilon \bs B_1$ into the Vlasov equation and eliminating all non-linear terms yields
\begin{equation}
    \dt f_0 + \bv \cdot \nabla_\bx f_0 + \frac{q}{m} \left(\bv \times \bs B_0\right) \cdot \nabla_\bv f_0 = \frac{q}{m v_\text{th}^2} f_0~ \bs E_1 \cdot \bv.
\end{equation}
We introduce the rescaled distribution function $g := f_1 / \sqrt{f_0}$. Maxwell's equations are linear by nature and remain unchanged. Dropping the 1 subscripts, the entire system writes
\begin{equation}
    \dt g_s + \bv \cdot \nabla_\bx g_s + \frac{q_s}{m_s} \left(\bv \times \bs B_0\right) \cdot \nabla_\bv g_s = \frac{q_s}{m_s v_\text{th,s}^2} \sqrt{f_{0,s}}~\bs E \cdot \bv
\end{equation}
\begin{align}
    \dt \bs D & = \curl \bs H - \bs J, \\
    \dt \bs B & = - \curl \bs E,
\end{align}
\begin{equation}
    \rho(t, \bx) = \sum_s q_s \int \sqrt{f_{0,s}}~ g_s(t, \bx, \bv) \dd\bv, \qquad \bs J(t,\bx) = \sum_s q_s \int \bv \sqrt{f_{0,s}}~ g_s(t, \bx, \bv) \dd\bv.
\end{equation}
With appropriate boundary conditions, the system conserves the energy
\begin{align}
    \cH & = \frac{1}{2} \int \bs E \cdot \bs D \dd\bx + \frac{1}{2} \int \bs B \cdot \bs H \dd\bx + \sum_s \frac{1}{2} m_s v_\text{th,s}^2 \int g_s^2 \dd\bx \dd\bv.
\end{align}
We apply a structure-preserving mimetic finite difference PIC discretization and obtain a discrete Hamiltonian system. In contrast to the non-linear case, the particle motion only depends on the background magnetic field but the particle weights are not constant anymore and interact with the electric field.
The discrete system and resulting splitting time integration scheme is given in \ref{sec:splitting_linvlasovmaxwell}.

Linearizing the equation has the advantage that only the perturbation of the distribution function has to be resolved. As a consequence, there is less PIC noise in the simulations and fewer particles are needed.

A standard test case is weak/linear Landau damping. The collisionless particle wave interaction phenomenon is started by a spatial variation in the particle density
\begin{equation}
    f(t=0,\bx,\bv) = (1 + \alpha \cos (\bk \cdot \bx)) f_0(\bx,\bv).
\end{equation}
The initial condition for the rescaled distribution function is
\begin{equation}
    g(t = 0, \bx, \bv) = \alpha \cos(\bk\cdot\bx) \sqrt{f_0 (\bx, \bv)}.
\end{equation}
The input file \verb|Examples/LinVlasovMaxwell/landauLVM.input| uses the parameters $k=0.4$ and $\alpha=0.04$. The resulting total energy composition and electric energy decay are shown in figures \ref*{fig:linvlasovmaxwell_energy} and \ref*{fig:linvlasovmaxwell_landau}, respectively.
\begin{figure}
    \centering
    \includegraphics[width=0.7\linewidth]{../_static/graphics/Examples/LinVlasovMaxwell_LandauTotalEnergy.png}
    \caption{Total energy composition during a Landau damping simulation.}
    \label{fig:linvlasovmaxwell_energy}
\end{figure}
\begin{figure}
    \centering
    \includegraphics[width=0.7\linewidth]{../_static/graphics/Examples/LinVlasovMaxwell_Landau.png}
    \caption{Landau damping: Exponential decay of the electric energy due to particle-wave interaction.}
    \label{fig:linvlasovmaxwell_landau}
\end{figure}

\end{document}